\chapter{Design}
After a manuscript has been written and marked up, it is time to create a visual
system that will emphasize the internal structure and the character of the
document. In print design, this involves the selection of one or several
typefaces that are well-suited to both the document and each other, the design
and the positioning of the structural elements of the document---such as
headings, tables, figures, and lists, and the choice of the paper size and the
page layout. In web design and multi-target publishing, several visual systems
may have to be created to accommodate for various display devices.

\section{Fonts}
When choosing \termpl{typeface} for a document, legibility should be of foremost
concern. The \term{body text} should be set with a typeface at a size of at least
10\,pt, if the document is aimed at adult readers, or 12\,pt, if visually
impaired readers and elementary-school students are a part of the
audience.~\cite[para.\,13--15]{kapr99}. The target medium also needs to be taken
into consideration. A faithful copy of a typeface designed for the letterpress
will look lighter than originally intended when printed digitally. This may
hamper its legibility, if it contains hairline
strokes~\cite[sec.~6.1.2]{bringhurst92}. In printed documents, typefaces with
\termpl{serif} are more familiar to the reader and therefore more suitable for
long-distance reading than their \term{sans serif} counterparts. At
low-resolution screens, however, simple low-contrast typefaces with slab or no
serifs will often yield the best result.

\begin{figure}
  \vspace{.5em}%
  \fbox{\includegraphics[clip,trim=4.3cm 19.8cm 4.35cm 4.55cm,%
    width=0.975\textwidth]{examples/03/cornford12.pdf}}%
  \vspace{.5em}%
  \inputcode[latex]{examples/03/cornford12.tex}
  \caption{An excerpt from \person{F. M. Cornford}'s \work{From Religion to
    Philosophy: A Study in the Origins of Western Speculation} as a
    text marked up in \TeX\index{TeX@\TeX} using \LaTeX{}\index{LaTeX@\LaTeX}
    macros and the primitives of \Hologo{XeTeX}\index{XeTeX@\Hologo{XeTeX}}
    (below) and the output document (above). Note that two typefaces were used:
    the regular typeface of Alegreya Sans at the size of 11\,pt for the Latin
    characters and the regular typeface of GFS Neohellenic at the size of
    13.2\,pt for the Greek characters.}
  \label{fig:cornford12}
\end{figure}

\begin{figure}
  \inputcode[html]{examples/03/cornford12.html}
  \caption{The document from Figure \ref{fig:cornford12} reformulated in
    \acronym{HTML}5 and \acronym{CSS}3}
\end{figure}

A typeface should also contain all the letters and symbols that will appear in
the document. If the manuscript is multilingual and contains passages in both
Latin and non-Latin writing systems, it may be necessary to combine several
typefaces. If the multilingual manuscript only contains Latin characters, but
several accented characters are missing from the body text typeface, they may be
constructed by combining the body text typeface with diacritical marks from
another font family. If certain punctuation marks and other symbols are missing
from the body text typeface, they may likewise be borrowed from other font
families. The typefaces should be consonant in their spirit and structure,
unless the text would benefit from the
dissonance.~\cite[sec.\,5.1.2]{bringhurst92}

Beside the body text typeface, several other typefaces may appear in a
document---a bold face, an italic face, or perhaps several sizes of the body
text typeface for use in the structural elements. The natural instinct is to
pick these typefaces from a single font family, but some families may not offer
all typefaces that the design requires. In these cases, the typefaces may again
have to be borrowed from other font families.

\section{Structural Elements}
\subsection{Paragraphs and Stanzas}\label{sec:paragraphs}
As the base units of linguistic thought in prose, \termpl{paragraph} split the
text into coherent portions ready for consumption. A line in a paragraph of the
body text should be 45--75 characters long on a single-column page or 40--50
characters long on a multi-column page and
\term*{justified}\index{alignment!justified}\index[see{alignment}]{justification}
(spread horizontally to fit the column width). Extended passages of lines wider
than 80 characters strain the eye of the reader, whereas justified lines that
are too narrow to accommodate 40 characters may make the word spacing entirely
too loose. In the latter case, the text should be set
\term*{ragged}\index{alignment!ragged}\index[see{alignment}]{rag} instead, as
seen in the sidenotes throughout this book~\cite[sec.\,2.1.2]{bringhurst92}.

% LaTeX, Word and the justification of paragraphs
% HTML, CSS and the justification and hyphenation of paragraphs

Vertically, the lines of a paragraph should be separated by approximately twenty
to forty-five percent of the typeface size~\cite{line-spacing}. If the size of
the body text typeface is 10\,pt, then the body text \term{line height} (also
known as the \term*{leading}\index[see{line spacing}]{leading}) would be between
12 and 14.5\,pt, adding 1 to 2.25\,pt of lead above and below each line. As a
general guideline, dark and bulky typefaces require more leading, as do texts
riddled with accents, full capital letters, subscripts, and
superscripts~\cite[sec.\,2.2.1]{bringhurst92}. The body text of this book is set
in 10\,pt Palatino with the leading of 12\,pt. To allow for such minimal
leading, all acronyms and other strings of upper-case letters are set as
\term{small capitals} (capital letters whose height matches the lower case).

% Leading in LaTeX, Word, and CSS

Two adjacent paragraphs should be visibly separated without distracting the
reader from the text. A predominant method is to indent the initial line of a
paragraph with one half (1\,en) to three times (3\,em) the typeface
size.\index{paragraph!indented} The indent is unnecessary when there is no
ambiguity---such as in the first paragraph following a
heading.~\cite[sec.\,2.3]{bringhurst92}

\noindent\hskip-.95em If the margins are ample, outdented paragraphs are an
intriguing option as well.\index{paragraph!outdented} \ding{161} Paragraphs can
also be separated by graphical symbols, such as pilcrows, bullets, or
boxes.\hspace{3em}A plain horizontal space that is at least 3\,em wide can
likewise act as a paragraph separator.~\cite[ch.\,2,~p.\,16]{beran14}

\noindent Block paragraphs\index{paragraphs!block} exchange indentation and
horizontal separators for additional vertical space above and below the
paragraph. In justified block paragraphs, this space can be omitted as well,
although the typesetter then has to manually ensure that the last line of each
paragraph offers enough horizontal space to act as a separator. In short
documents and limited spans of text, block paragraphs are an attractive
option.~\cite[sec.\,2.3.2]{bringhurst92}

% Paragraph separation in LaTeX, Word, and CSS

Being the verse counterpart to the paragraph, the stanza is a collection of
lines rather than of sentences. Due to this structural difference, stanzas are
typically only justified, when the individual lines are long enough to fill up
the column, and ragged otherwise. Much like in the case of prose, short-form
poetry benefits from having the stanzas set in block paragraph style.

\subsection{Headings}
Another fundamental structural element is the \term{heading}. The function of a
heading is to delimit and name the individual sections of a document. To
alleviate navigation, headings should be a prominent presence on a page. This
can be achieved by using a larger variant of the body text typeface or by
including the text of the latest heading in the margin or the header of the
page~\cite[sec.\,4.2.1]{bringhurst92}, as seen throughout this book.

The hierarchy of the headings can be expressed through the variation of
typefaces, indentation, alignment and numbering, although alternating the size
of the body text typeface is sufficient for many types of documents. In
documents that are bound in codex form and read two pages at a time, the height
of headings should be a whole multiple of the line height of the body text, so
that the headings do not disrupt the alignment of lines on the facing
pages.~\cite[para.\,33]{kapr99}

% Headings in LaTeX, Word and CSS

\subsection{Tables and Lists}
\Termpl{table} and \termpl{list} are structural elements that should fit
seamlessly into the surrounding text and avoid unnecessary visual clutter. Use
the same typeface the surrounding text does, treat the columns of tables the
same way you treat columns in the text, and keep the amount of rules, boxes,
dots, and extraneous spacing to a bare minimum (see
Table~\ref{table:papers}).~\cite[sec.\,2.1.10 and 4.4]{bringhurst92}

\begin{table}
  \begin{center}
    \hspace{4ex}Sizes in inches\hspace{1.2em}Page proportions
    \begin{tabular}{%
        l%
        D{×}{{}\times{}}{4}@{\quad\quad}%
        D{:}{{}:{}}{5}@{\quad\quad}%
        D{.}{.}{5}}
      \emph{A4}     & 8.27×11.7 & 2:\sqrt{2} & 1.41421 \\
      \emph{B5}     & 6.93×9.84 & 1:\sqrt{2} & 0.707   \\
      \emph{Letter} &   8½×11   & 1:1.294    & 1.2941  \\
    \end{tabular}
  \end{center}
  \caption{An overview of common paper sizes used for commercial and industrial
    printing}
  \label{table:papers}
\end{table}

\subsection{Notes}
\Termpl{note} provide commentary on a specified passage of the main text and can
take three different forms:
\begin{enumerate}
  \item[1]\Termpl{sidenote} are displayed in the horizontal margins next to the
    \sidenote{%
      This is a sidenote. Sidenotes enliven the page and are easy for the reader
      to find.}
    relevant passage of the main text, as seen throughout this book.  Unless the
    horizontal margins are very wide, sidenotes are unsuitable for the inclusion
    of bibliographical references---a common use for notes in academic writing.
  \item[2]\Termpl{footnote} are delegated to the bottom of the page and linked
    to the relevant passage of the main text through symbols or superscript
    numbers.\kern-.3ex\footnote{%
      This is a footnote. Due to their width, footnotes can
      comfortably accommodate full bibliographical references, which makes them
      popular in academic writing.
      
      A footnote can also contain multiple paragraphs of text, although long
      footnotes are tedious to read if the size of the typeface is
      small.~\cite[sec.\,4.3.1]{bringhurst92}}
    Compared to side notes, they are more difficult for the reader to find.
    Footnotes should align with the bottom of the text block, not stick out into
    the bottom margin.~\cite[para.\,48]{kapr99}
  \item[3]\Termpl{endnote} are delegated to the end of a section or the entire
    document and are linked to the relevant passage of the body text through
    superscript numbers. They are the easiest of the three to typeset, but also
    the hardest for the reader to find.
\end{enumerate}

Notes are typically typeset in sizes from 8\,pt up to the body text typeface
size depending on their frequency, importance, and average
length.~\cite[sec.\,4.3]{bringhurst92} If several categories of notes are
present in the document, it may be desirable to give each a different form.

\subsection{Quotations}
Quotations repeat what has already been expressed somewhere else before and can
take two different forms:~\cite[sec.\,5.4]{bringhurst92}
\begin{enumerate}
  \item[1]\Termpl*{run-in quotation}\index{quotation!run-in}
    are included directly into the paragraph and set off from the surrounding
    text using quotation marks in accordance with the orthographic rules on the
    use of punctuation in the language of the paragraph: \quote{Jesters do oft
    prove prophets.}
    
    From the designer's viewpoint, run-in quotations require no special
    treatment, although it is crucial that the body text typeface contains the
    required quotation marks.
  \item[2]\Termpl*{block quotation}\index{quotation!block}
    are set as block paragraphs\index{paragraph!block} that are clearly
    separated from the surrounding text. This involves adding a vertical space
    above and below the block paragraphs and optionally also changing the
    typeface, its size, or the indentation of the
    paragraphs:~\cite[sec.\,2.3.3]{bringhurst92}\blockquote{%
      This is the excellent foppery of the world that when we are sick in
      fortune—often the surfeit of our own behavior—we make guilty of our
      disasters the sun, the moon, and the stars, as if we were villains by
      necessity, fools by heavenly compulsion, knaves, thieves, and treachers by
      spherical predominance, drunkards, liars, and adulterers by an enforced
      obedience of planetary influence, and all that we are evil in by a divine
      thrusting-on. An admirable evasion of whoremaster man, to lay his goatish
      disposition to the charge of a star!}{\person{William Shakespeare},
      \work{King Lear}}
    Block quotations are ideal for longer quotations and for quotations that
    should carry more weight that run-in quotations.
\end{enumerate}

\section{Page Layout}
The page consists of a textblock surrounded by margins. The text width
area is largely determined by the number of columns and the body text size---as
described in Section~\ref{sec:paragraphs}---as well as by our plans for the
horizontal margins. A margin containing an occasional sidenote will require less
space that a margin ripe with photographs, tables, and diagrams.

The vertical margins may contain additional navigational aids, such as the page
numbers and running headers in this book.~\cite[sec.\,8.5.2]{bringhurst92}

In print design---and wherever else the page height is fixed---we need to also
decide on the text height. The text height needs to be a multiple of the body
text line height, so that it is possible to completely fill the text block with
text.  It is typical to derive the text height from the text width to achieve
proportions that work well with the proportions of the
page.~\cite[sec.\,8.4.2]{bringhurst92}

\section{Color}
Three types of \termpl{cone cell}---S, M, and L---exist in the human eye, and
each type is sensitive to a different range of the \term{visible light}.  The
perception of the visible light is what is referred to as \term{color}. In this
section, we will briefly overview the color theory as well as general
guidelines with respect to color scheme design.

\begin{figure}
  \includegraphics[width=0.5\textwidth]{examples/03/rgb}%
  \includegraphics[width=0.5\textwidth]{examples/03/cmy}
  \caption{The additive $RGB$ color model (left), and the subtractive $CMY$
    color model (right).}
  \label{fig:color-models}
\end{figure}

\subsection{Theory}
In \citeyear{cie32}, \acronym{CIE} defined the \term*{\acronym{CIE} $XYZ$ color
space}\index{color space}\index{color space!$XYZ$}, which specifies how a
\acronym{SPD} of the visible light relates to a set of three parameters ($X,
Y,$ and $Z$) that specify a color. Since the parameter $Y$ is a measure of the
\term{luminance} of a color, a two-dimensional color space called the
\term*{\acronym{CIE} $xyY$ color space}\index{color space!$xyY$} can by derived
from \acronym{CIE} $XYZ$ by disregarding the luminosity.
A plot, where a point $(x, y)$ corresponds to a color in \acronym{CIE} $xyY$,
is called a \term{chromacity diagram} (see Figure~\ref{fig:chromacity-diagram}).
In theory, we can directly use the \acronym{CIE} color spaces to specify
colors. In practice, there exist task-specific color spaces.

\begin{figure}
  \includegraphics[width=\textwidth]{examples/03/CIExy1931}
  \caption{A \acronym{CIE} $xyY$ chromacity diagram, where the shark-fin-shaped
    \term{spectral locus} corresponds to a narrow-band \acronym{SPD} with power
    at just a single wavelength swept across the visible light (380–700\,nm).}
  \label{fig:chromacity-diagram}
\end{figure}

\begin{figure}[b!]
  \includegraphics[width=\textwidth]{examples/03/1500-latin-vulgate}
  \caption{An excerpt from the \work{Latin Vulgate Bible} printed by
    the German goldsmith, printer, and publisher \person{Anton Koberger} in
    1487}
  \label{fig:vulgate-bible}
\end{figure}

Electronic screens use red, green, and blue sub-pixels. This gives rise to the
additive \term*{$RGB$ color model}\index{color model!$RGB$}
(see Figure~\ref{fig:color-models}). An \term*{$RGB$ color space}\index{color
space!$RGB$} is then specified by the chromacities of its
\termpl{primary color} (\termpl*[primaries]{primary}), and by the \term{white
point}. The extent of the colors that can be represented by an $RGB$ color
space (its \term{gamut}) can be plotted in the chromacity diagram as a triangle
whose vertices are the chromacities of the primaries. Standard $RGB$ color
spaces include sRGB, Adobe RGB 1998, and ProPhoto RGB (see
Figure~\ref{fig:chromacity-diagram})\index{color space!sRGB}\index{color
space!Adobe RGB 1998}\index{color space!ProPhoto RGB}.
There exists no single standard white point or primaries; if only the $RGB$
coordinates of a color are given, the color cannot be accurately reproduced.
For this reason, digital images often include an \acroshort{ICC} \term{color
profile}, which fully specifies a color space.

Mixing cyan, magenta, and yellow inks gives rise to the subtractive
\term*{$CMYK$ color model}\index{color model!$CMYK$}. Although three colors ($C,
M,$ and $Y$) are sufficient to describe a color, a black ink ($K$) is less
expensive and dries more quickly than a combination of three colored inks.
White color is provided by the paper. Similarly to the $RGB$ model, $CMYK$
coordinates specify a color only if we specify a $CMYK$ color space\index{color
space!$CMYK$} using an \acroshort{ICC} color profile. For printing text, it is
often advantageous \emph{not} to specify a color profile, so that no inks other
than $K$ are introduced in a \term{color conversion}. When an accurate
reproduction of a color is desired, a color model must be specified.

\subsection{Schemes}
In both print and web design, it is perfectly reasonable to use either just the
combination of black and white or shades of gray. A secondary color may be
introduced to enliven the page if necessary; \textcolor{red}{\textbf{red}} has
historically been used for this purpose (see Figure~\ref{fig:vulgate-bible}).
More than a single hue of color may be introduced, although each additional one
makes it more difficult to establish a visual system that is intelligible to
the reader.

You should only use colored typefaces for emphasis, not for the body text, and
on backgrounds that are white or of sufficient contrast with the typeface
color. Distinct colors should stay distinct even for the color-blind reader,
unless the lack of distinction between the colors does not impair
understanding.
